\begin{sol}{xca:neutral}
If $e$ and $e_1$ are both neutral elements, then $e=ee_1=e_1$. 
\end{sol}

\begin{sol}{xca:ax=b}
    If $ax=b$, after multiplying on the left by $a^{-1}$ we
    obtain that $x=a^{-1}b$. Similarly, the equation $xa=b$ 
    has $x=ba^{-1}$ as its unique solution. 
\end{sol}

\begin{sol}{xca:LR}
    For $g\in G$, the map $L_g\colon G\to G$, $x\mapsto gx$, is invertible 
    with inverse $L_{g^{-1}}$, as
    \[
    (L_g\circ L_{g^{-1}})(x)=g(g^{-1}x)=(gg^{-1})x=x
    \]
    for all $x\in G$. Similarly, $L_{g^{-1}}\circ L_g)(x)=x$ for all $x\in G$. 
    
    In the same way, we prove that 
    for each $g\in G$, the map $R_{g^{-1}}$ is the inverse of $R_g$. 
\end{sol}

\begin{sol}{xca:GxH}
    To prove the associativity, let $g,g_1,g_2\in G$ and 
    $h,h_1,h_2\in H$. Since $G$ and $H$ are groups, 
    their multiplications are associative. Then 
    \begin{align*}
        ((g,h)(g_1,h_1))(g_2,h_2) &= (gg_1,hh_1)(g_2,h_2)\\
        &=((gg_1)g_2,(hh_1)h_2)\\
        &= (g(g_1g_2),h(h_1h_2))\\
        &= (g,h)(g_1g_2,h_1h_2)\\
        &= (g,h)((g_1,h_1)(g_2,h_2)).
    \end{align*}
    
    The neutral element of $G\times H$ is $(1,1)$, as $(1,1)(g,h)=(g,h)=(g,h)(1,1)$. 
    
    The inverse
    of $(g,h)$ is $(g,h)^{-1}=(g^{-1},h^{-1})$, as
    \begin{align*}
    (g,h)(g,h)^{-1}&=(g,h)(g^{-1},h^{-1})=(gg^{-1},hh^{-1})=(1,1),\\
    (g,h)^{-1}(g,h)&=(g^{-1},h^{-1})(g,h)=(g^{-1}g,h^{-1}h)=(1,1).
    \end{align*}
\end{sol}


